\scanpage{59}
\mspage{46}
% archive sheet number: 58

\[
 J'=J\cup\{i_0\}=J\cup I_{/i_0}
\]

et montrons qu'on peut prolonger le \deleted{\unsure{quadru}} \inserted{triplet} $Z,i,p$ en $Z',i',p'$ sur $J'$. Posant $\deleted{\unsure{I}}K'=I_{/i_0}$, $K=J\cap K$\inserted{$=I_{<i_0}$}, on a donc

\diagram{\begin{tikzcd}[ampersand replacement=\&,column sep=large]
 K=J\cap K'\arrow[r,hook]\arrow[d,hook] \& K'=I_{/i_0}\arrow[d,hook]\\
 J\arrow[r,hook] \& J'=J\cup K'
\end{tikzcd}}

On voit donc que les prolongements de $(Z,i,p)$ de $J$ à $J'$ correspondent à ceux de $(Z|K',i|K',p|K')$ à $K$. Pour ceci, on note que pour une flèche $g$ \deleted{\unsure{d'objets}} objets de $M(J'^{\unsure{\circ|0}})$, cette flèche est dans $C_{J'^\circ}$ resp. $W_{J'^\circ}$, ssi ses restrictions à $J$ et à $K'$ le sont.

On est donc ramené au cas \unsure{standart|standard} d'un ensemble ordonné \deleted{\unsure{$K$ ayant}} avec objet final $\unsure{e_0|i_0}$, ens. ordonné qu'on va noter $I$, et l'ens ordonné $I'$ des objets $i<e_0$, i.e. $I'=I\setminus\{e_0\}$, on a $f:X_{\unsure{\cdot|\bullet}}\to Y_{\cdot}$ \deleted{\unsure{flèche}} dans $M(I^\circ)$, une factorisation de $f|I'$ en $p'\circ i'$ avec $i'\in C_{I'^\circ}$, $p'\in W_{I'^\circ}$,
